\documentclass[11pt]{article}

\usepackage{amsmath}
\usepackage{amssymb}
\usepackage{geometry}
\usepackage{listings}
\usepackage{xcolor}

\geometry{margin=2.5cm}

% ------------------------------------------------------------
% Julia listings
% ------------------------------------------------------------
\lstdefinelanguage{Julia}{
  keywords={function,end,for,if,else,return,using,struct},
  sensitive=true,
  comment=[l]#,
  morestring=[b]"
}

\lstset{
  language=Julia,
  basicstyle=\ttfamily\small,
  keywordstyle=\color{blue},
  commentstyle=\color{gray},
  frame=single,
  breaklines=true
}

\title{Validation Tests for \texttt{avoid_points_path}}
\author{}
\date{}

\begin{document}

\maketitle

\section{Introduction}

This document collects validation tests for the function
\texttt{avoid_points_path}, designed to construct discrete polygonal
paths in the complex plane avoiding a finite set of forbidden points.
The tests verify correctness, geometric safety, smoothness, and numerical
stability, which are essential for reliable path lifting in
\texttt{PolynomialCovers.jl}.

\section{Auxiliary Metrics}

\subsection{Minimum Distance to Forbidden Set}

\begin{lstlisting}
function min_distance_to_set(p::Path, S)
    minimum(
        abs(v.z - a.z)
        for v in p.vertices
        for a in S
    )
end
\end{lstlisting}

\subsection{Maximum Turning Angle}

\begin{lstlisting}
function max_turn_angle(p::Path)
    θmax = 0.0

    for i in 2:length(p.vertices)-1
        v1 = p.vertices[i].z - p.vertices[i-1].z
        v2 = p.vertices[i+1].z - p.vertices[i].z

        θ = acos(
            real(conj(v1)*v2) /
            (abs(v1)*abs(v2))
        )

        θmax = max(θmax, θ)
    end

    θmax
end
\end{lstlisting}

\section{Tests}

\subsection{Basic Sanity Test}

\begin{lstlisting}
@testset "avoid_points_path basic" begin

    z0 = Vertex(0.0,0.0)
    zf = Vertex(5.0,0.0)

    S = [
        Vertex(2.5,0.5),
        Vertex(2.5,-0.5)
    ]

    p = avoid_points_path(z0,zf,S)

    @test startpoint(p) == z0.z
    @test endpoint(p)   == zf.z
    @test length(p) > 2
end
\end{lstlisting}

\subsection{Safety Distance Test}

\begin{lstlisting}
@testset "safety distance" begin

    z0 = Vertex(0.0,0.0)
    zf = Vertex(6.0,0.0)

    S = [Vertex(3.0,0.0)]

    p = avoid_points_path(z0,zf,S)

    dmin = min_distance_to_set(p,S)

    @test dmin > 0.2
end
\end{lstlisting}

\subsection{Obstacle Bypass}

\begin{lstlisting}
@testset "obstacle bypass" begin

    z0 = Vertex(0.0,0.0)
    zf = Vertex(4.0,0.0)

    S = [Vertex(2.0,0.0)]

    p = avoid_points_path(z0,zf,S)

    ys = [imag(v.z) for v in p.vertices]

    @test maximum(abs.(ys)) > 0.1
end
\end{lstlisting}

\subsection{Smoothness Test}

\begin{lstlisting}
@testset "smoothness" begin

    p = avoid_points_path(
        Vertex(0,0),
        Vertex(5,0),
        [Vertex(2.5,0.0)]
    )

    @test max_turn_angle(p) < pi/2
end
\end{lstlisting}

\subsection{Allowed Final Collision}

\begin{lstlisting}
@testset "final allowed collision" begin

    z0 = Vertex(0,0)
    zf = Vertex(3,0)

    S = [zf]

    p = avoid_points_path(z0,zf,S)

    hits =
        count(v -> v == zf, p.vertices)

    @test hits == 1
end
\end{lstlisting}

\subsection{Random Stability Test}

\begin{lstlisting}
@testset "random stability" begin

for k in 1:20

    z0 = Vertex(0.0,0.0)
    zf = Vertex(5.0,0.0)

    S = [Vertex(rand()*4, rand()-0.5)
         for _ in 1:5]

    p = avoid_points_path(z0,zf,S)

    @test endpoint(p) == zf.z
end

end
\end{lstlisting}

\section{Conclusion}

These tests validate that the generated paths:

\begin{itemize}
\item reach the prescribed endpoint,
\item avoid forbidden points,
\item maintain safety distance,
\item remain geometrically smooth,
\item behave robustly under random configurations.
\end{itemize}

Such properties are essential for stable numerical lifting procedures
in polynomial covering algorithms.

\end{document}